
\section{Statik}
\paragraph{Rechte-Hand-Regel}
$\vec{a}$-Daumen, $\vec{b}$-Zeigefinger $\to$ $\vec{a}\times\vec{b}$-Mittelfinger
\paragraph{Drehmoment}\
\begin{definition}
	$M=r\times F$
\end{definition}
\paragraph{Gleichgewichte in der Statik}
\begin{itemize}
	\item Kräftegleichgewicht
	\item Momentgleichgewicht
\end{itemize}
Freischnitt
\paragraph{Arten von Lagern}
\begin{center}
	\setlength{\extrarowheight}{4pt} 
	\begin{tabularx}{0.8\textwidth}{>{\bfseries\arraybackslash}l|X|X|X}
		{\large Lagerwertigkeit}&einwertig&zweiwertig&dreiwertig\\
		\hline
		Beispiel&verschiebbare Gelenklager, verschiebbare Parallelführung& festes Gelenklager, Parallelführung& starre Einspannung
\end{tabularx}
\end{center}
\paragraph{Statisches Gleichgewicht}
reduziert im $\R^2$ auf 3 Gleichungen:
\begin{align*}
	F_x=0, F_y=0, M_Z=0
\end{align*}
\paragraph{Statische Bestimmtheit}
Drei Gleichungen:
\begin{itemize}
	\item[$\to$] mind. drei Unbekannte, sonst unterbestimmt (verschieblich)
	\item[$\to$] max. drei Unbekannte, sonst überbestimmt (unbestimmt)
\end{itemize}
Also:
\begin{itemize}
	\item exakt drei Unbekannte (notwendige Bedingung)
	\item eindeutig lösbar (lineare Unabhängigkeit, Determinante der Koeffizientenmatrix $\neq 0$)
\end{itemize}
\begin{align*}
	A\vec{V}=\vec{F}
\end{align*}
\begin{center}
	\begin{tabularx}{0.9\textwidth}{XXX}
	{\bfseries Oberflächenkräfte}&{\bfseries Volumenkräfte}&{\bfseries Streckenlast}\\
	$\begin{aligned}
	\vec{F}&=\int\limits_{(S)}\vec{f}_S\odif{S}\\
	\vec{M}_A&=\int\limits_{(S)}(\vec{r}_{AP}\times\vec{f}_S)\odif{S}
\end{aligned}$&$\begin{aligned}
	\vec{F}&=\int\limits_{(V)}\vec{f}_V\odif{V}\\
	\vec{M}_A&=\int\limits_{(V)}(\vec{r}_{AP}\times\vec{f}_V)\odif{V}
\end{aligned}$&
$\begin{aligned}
	F &= \int\limits_0^l q(\xi)\odif{\xi}\\
	M &= \int\limits_0^l \xi\cdot q(\xi)\odif{\xi}
\end{aligned}$
\end{tabularx}
\end{center}


\section{Festigkeitslehre}

\section{Dynamik}

